\section{Finite ratings and estimation error}
\subsection{Noise in the observed target mean}
Let $I$ be a random held-out item and $Y_{I1},\ldots,Y_{Ik}$ be conditionally independent ratings drawn from $P_I$.
Set $t_I=\E[Y\mid I]$, $\sigma_I^2=\Var(Y\mid I)$, and $\widehat t_I=k^{-1}\sum_{j=1}^kY_{Ij}=t_I+\varepsilon_I$.
Then
\[
 \E[\varepsilon_I\mid I]=0,\qquad \Var(\varepsilon_I\mid I)=\frac{\sigma_I^2}{k}.
\]
For a fixed forecast mean $m_I$, or a forecast conditioned on training data independent of these test ratings, the law of total covariance yields
\[
 \Cov(m_I,\varepsilon_I)=\E[(m_I-\E m_I)\E(\varepsilon_I\mid I)]=0.
\]
The same argument gives $\Cov(t_I,\varepsilon_I)=0$.
The law of total variance therefore implies
\[
 \Cov(m_I,\widehat t_I)=\Cov(m_I,t_I),\qquad
 \Var(\widehat t_I)=\Var(t_I)+\E[\sigma_I^2]/k.
\]
Whenever the relevant variances are positive,
\begin{equation}\label{eq:attenuation}
 \operatorname{Corr}(m_I,\widehat t_I)
 =\frac{\operatorname{Corr}(m_I,t_I)}{\sqrt{1+\E[\sigma_I^2]/(k\Var(t_I))}}.
\end{equation}
Conditional mean-zero noise suffices; independence of the noise from the item is unnecessary and generally false with heteroskedastic ratings.
For item-dependent counts $k_I$, replace $\E[\sigma_I^2]/k$ by $\E[\sigma_I^2/k_I]$ under the same sampling assumptions.

Equation~\eqref{eq:attenuation} is a population correlation identity over random held-out items and fresh ratings.
It is not an equality for the expectation of a finite-sample Pearson coefficient, whose denominator is random.
It also does not apply directly to a predictor fitted on the same rating noise.
An unrestricted predictor can interpolate $m_i=a+b\widehat t_i$ whenever feasible and the observed means are nonconstant, attaining training correlation one while memorizing noise.
Correlated raters, changing rater populations, and selection of which items receive ratings require a different noise analysis.

\subsection{A regular asymptotic CE calculation}
The familiar $P/(2N)$ expression concerns a specific estimation regime, not every distributional learner.
To state it precisely, let $Z_j=(X_j,Y_j)$, $j=1,\ldots,N$, be independent and identically distributed observations from a well-specified parametric conditional model $q_{\theta^\ast}(Y\mid X)$.
Let $\theta^\ast$ lie in the interior of an identifiable $P$-dimensional parameter domain.
Assume a twice differentiable population CE risk
\[
 R(\theta)=\E[-\log q_\theta(Y\mid X)]
\]
with nonsingular Fisher information $I=\nabla^2R(\theta^\ast)$, and an efficient consistent unpenalized maximum-likelihood estimator satisfying
\[
 \sqrt N(\widehat\theta-\theta^\ast)\ \xrightarrow{d}\ \mathcal N(0,I^{-1}).
\]
Assume in addition the moment convergence needed to take expectations of the quadratic form and an expected Taylor remainder $o(N^{-1})$.
For example, appropriate uniform integrability of $N\|\widehat\theta-\theta^\ast\|^2$ together with control of the remainder suffices.
Then
\begin{align}
 R(\widehat\theta)-R(\theta^\ast)
 &=\tfrac12(\widehat\theta-\theta^\ast)^\top I(\widehat\theta-\theta^\ast)+\text{remainder},\\
 \E[R(\widehat\theta)-R(\theta^\ast)]
 &=\frac1{2N}\operatorname{tr}(II^{-1})+o(N^{-1})
 =\frac{P}{2N}+o(N^{-1}).\label{eq:asymptotic}
\end{align}
The first-order Taylor term vanishes at the interior population minimizer.
The second equality uses $\E[N(\widehat\theta-\theta^\ast)(\widehat\theta-\theta^\ast)^\top]\to I^{-1}$ and the trace identity for quadratic forms.
Thus the expectation statement requires more than convergence in distribution alone.
The result is a conditional regular-asymptotic calculation, measured in nats per fresh observation.

When $N$ is comparable to $P$, when the MLE lies near a boundary or does not exist, under misspecification, or with strong regularization, these assumptions can fail.
Overparameterized models and clustered ratings need separate analyses.
In particular, an estimated zero probability can yield infinite test CE even when the corresponding category was absent from training.
A probability floor changes the estimator and its evaluated loss; it is not implicit in the preceding identities.

\subsection{Why there is no universal break-even sample size}
Correlation's population optimum set includes the truth when the class is well specified.
Its nonuniqueness therefore does not force a positive error floor for every fitting procedure.
Initialization, regularization, a tie rule, and model restrictions can select different members of that set.
If a specified procedure converges to a forecast with positive excess CE $B$, then comparing $P/(2N)$ to $B$ suggests the heuristic scale $N\approx P/(2B)$.
This comparison requires both risks to use CE in nats and the asymptotic assumptions to be relevant.
It is invalid to substitute an RPS gap for $B$ in that expression.
Neither $B>0$ nor a universal crossover follows from correlation's invariances.

Properness describes the optimizer of population expected loss.
It does not order finite-sample estimators, optimization methods, or validation procedures.
With a small number of ratings, a restricted or regularized forecast can have lower held-out RPS or CE than a more flexible fit.
The invariance results also do not compare choosing among existing judges with fitting a new distributional head.
These are distinct statistical procedures whose performance depends on the candidate class and data.
